
\providecommand{\ClaimStatusConjectured}{}
\providecommand{\ClaimStatusHeuristic}{}
\providecommand{\ClaimStatusConditional}{}
\providecommand{\ClaimStatusProvedElsewhere}{}
\providecommand{\ClaimStatusOpen}{}
\providecommand{\ClaimStatusFolklore}{}

\chapter{The mirror discriminant}
\label{ch:view4-mirror}

\section{The mirror discriminant}
\label{sec:view4-mirror-discriminant}

The automorphic, BKM, Pfaffian, and theta-lift constructions place
\(\Delta_5\) on chains that begin with explicit modular or
representation-theoretic data: the
Borcherds--Gritsenko--Nikulin section of the weight-\(5\) automorphic
line over \(\Sp_4(\Z) \backslash \mathcal H_2\) (View~\textcircled{1}),
the denominator
\(\operatorname{den}(\mathfrak g_{\Delta_5}) = 64^{-1} \Delta_5(2Z)\) of
the Borcherds--Kac--Moody Lie superalgebra on \(\Lambda_{II}^{2,1}\)
(View~\textcircled{2}), the protected Pfaffian
\(\operatorname{Pf}_{\mathrm{prot}}(\mathfrak D_X^{\mathrm{DI}})\) of the
Dirac--Igusa pro-object (View~\textcircled{3}), and the singular theta
lift of the weak Jacobi form \(\phi_{0,1}\), namely the Borcherds
regularized theta integral on the
\(\mathrm O(3,2)\) type-IV domain, with \(\Mp_2\)-Weil input and the
\(\PSp_4\simeq\SO(3,2)_+\) exterior-square bridge
(View~\textcircled{5}).
The period-side assertion is geometric. Let \(\mathcal P_{K3\times E}^{(3)}\) be
the rank-\(3\) Mukai-invariant period quotient transported to
\(\mathcal H_2/\Sp_4(\Z)\).  Let \(\mathfrak H_{\mathrm{II}}\) denote
the type-II Heegner divisor on this quotient.  The automorphic
calculation gives the divisor equality
\begin{equation}
\label{eq:view4-conjecture}
[\,\Delta_5^{-2}\,]
\;=\;
2[\mathfrak H_{\mathrm{II}}]
\qquad
\text{(equality of divisor classes on the }
(K3 \times E)\text{-component)},
\end{equation}
where \([\Delta_5^{-2}]\) is the polar divisor of the meromorphic
section \(\Delta_5^{-2}\) of the automorphic line.  The conjectural
content is the existence of a rank-\(3\) Mukai--Hodge quotient of a
mirror period family whose discriminant support is
\(\mathfrak H_{\mathrm{II}}\), together with the comparison of the
automorphic line with the period Hodge line.

This identification is conjectural. The cross-level identities
\textcircled{4}\(\leftrightarrow\)\textcircled{1},
\textcircled{4}\(\leftrightarrow\)\textcircled{2},
\textcircled{4}\(\leftrightarrow\)\textcircled{5} are open.  The
automorphic identity, the BKM denominator identity, the conditional
Pfaffian identity, and the singular theta-lift identity do not use the
period-discriminant conjecture.

The handle \(\Delta_5\) splits into three names. The automorphic section over
\(\Sp_4(\Z) \backslash \mathcal H_2\) is denoted
\(\mathcal D_X = \Delta_5\) (View~\textcircled{1}); the BKM denominator
is
\(\operatorname{den}(\mathfrak g_{\Delta_5}) = 64^{-1} \Delta_5(2Z)\)
(View~\textcircled{2}); the protected Pfaffian is
\(\operatorname{Pf}_{\mathrm{prot}}(\mathfrak D_X^{\mathrm{DI}})\)
(View~\textcircled{3}). The mirror conjecture concerns the first: the
divisor identification~\eqref{eq:view4-conjecture} is a statement about
\(\mathcal D_X\) as an automorphic section, transported to the period
side.

\subsection{The K3 mirror-period family}
\label{subsec:k3-mirror-period}

The lattice-polarised K3 mirror is fixed by Dolgachev~\cite{DolgachevMirrorLP}
and tracked, throughout, through the Mukai lattice. Let
\[
\Lambda_{K3}
\;=\;
3U \oplus 2 E_8(-1),
\qquad
\widetilde \Lambda_{K3}
\;=\;
U \oplus \Lambda_{K3}
\;\simeq\;
4U \oplus 2 E_8(-1),
\]
the K3 lattice of signature \((3,19)\) and the Mukai lattice of
signature \((4,20)\) (Mukai~\cite{Mukai1984,Mukai1987}). A
\emph{lattice polarisation} of a complex K3 surface \(S\) is a
primitive embedding \(\iota \colon M \hookrightarrow \Pic(S)\) of an
even hyperbolic lattice \(M\), of signature \((1, t)\), such that
\(M^\perp \subset \Lambda_{K3}\) admits a primitive embedding of the
hyperbolic plane \(U\); the pair \((S, \iota)\) is then
\emph{Nikulin-admissible}, in the convention of
Gritsenko--Nikulin~\cite{GNII}\,\S\,3 and the rank-3 sublattice
constructions of Borcherds~\cite{Borcherds95}\,\S\,15. Here
\[
t=\rk M-1
\]
is the negative index of \(M\), so that
\(\rk M^\perp=21-t\) and
\(\mathrm{sign}(M^\perp)=(2,19-t)\).  The convention is
\(\mathrm{sign}(L)=(\#\text{positive},\#\text{negative})\).

The associated period domain
\[
\Omega_M
\;=\;
\bigl\{
[\omega] \in \P\!\bigl(M^\perp \otimes \C\bigr)
\;:\;
\langle \omega, \omega \rangle = 0,
\;
\langle \omega, \bar\omega\rangle > 0
\bigr\}^+
\]
is the type-IV symmetric domain of signature \((2, 19 - t)\) and
dimension \(19-t\), with arithmetic group
\(\Gamma_M = \mathrm O^+(M^\perp)\) acting properly discontinuously.
The \emph{Dolgachev mirror lattice} is the
sub-lattice
\[
\check M
\;=\;
M^\perp \cap (U)^\perp
\;\subset\;
\Lambda_{K3},
\qquad
M^\perp
\;=\;
U \oplus \check M,
\]
so that \(\check M\) carries signature \((1,18-t)\) and is hyperbolic.
The mirror K3 surface \(S^\vee\) is lattice-polarised by \(\check M\).
Dolgachev mirror symmetry gives a correspondence between the
\(M\)-polarised and \(\check M\)-polarised moduli quotients; it is an
involution on a single quotient only in the self-mirror cases.

The Aspinwall--Morrison form of the same construction~\cite{AspinwallMorrisonStringy}
uses the Mukai lattice to compare complex and complexified K\"ahler
directions at large-volume and large-complex-structure cusps.  The full
K3 sigma-model moduli space is larger than \(\Omega_M/\Gamma_M\);
the lattice correspondence is the only input needed for the
\(\Delta_5\) discriminant statement.  The type-IV lattice of the mirror
quotient is
\[
\widetilde M
\;=\;
U_{\mathrm{mir}} \oplus M
\;\simeq\;
\check M^\perp
\;\subset\;
\Lambda_{K3},
\qquad
\mathrm{sign}(\widetilde M)=(2,t+1).
\]
Here \(U_{\mathrm{mir}}\) is the hyperbolic plane split from
\(M^\perp=U_{\mathrm{mir}}\oplus\check M\).  In the full Mukai lattice
\(\widetilde\Lambda_{K3}=U_{\mathrm{Muk}}\oplus\Lambda_{K3}\), the
orthogonal complement of \(U_{\mathrm{Muk}}\oplus M\) is
\(M^\perp=U_{\mathrm{mir}}\oplus\check M\); the extra Mukai plane is
contained in the augmented sublattice.  No canonical action on
\(\widetilde \Lambda_{K3}\) is needed here; only the primitive
decomposition and the induced period domains enter the argument.

\paragraph{The rank-three case.} The polarisation relevant to View~\textcircled{4}
is
\[
M
\;=\;
U \oplus \langle -2 \rangle,
\qquad
\mathrm{rank}(M) = 3,
\qquad
\mathrm{sign}(M) = (1, 2),
\]
with Dolgachev mirror lattice
\[
\check M
\;=\;
\bigl(M^\perp \cap U^\perp\bigr)_{\Lambda_{K3}}
\;\simeq\;
U \oplus E_7(-1) \oplus E_8(-1),
\]
of rank \(17\) and signature \((1, 16)\) inside \(\Lambda_{K3}\); the
convention is
\(\mathrm{sign} = (\#\text{positive}, \#\text{negative})\) on the
bilinear form of the lattice. The decomposition follows from the
embedding \(\langle -2 \rangle \hookrightarrow E_8(-1)\) on a single
root, after which the orthogonal complement of the root inside
\(E_8(-1)\) is \(E_7(-1)\). The rank-five mirror-domain lattice is
\[
\widetilde M
\;=\;
U \oplus M
\;=\;
2 U \oplus \langle -2 \rangle,
\qquad
\mathrm{rank}(\widetilde M) = 5,
\qquad
\mathrm{sign}(\widetilde M) = (2, 3).
\]
Under the exterior-square convention (the model of
\(\PSp_4(\R)\simeq\SO(3,2)_+\) at \S\,\ref{ssec:exceptional-iso})
the same lattice is written
\(\Lambda^{3,2} = U \oplus U \oplus \langle 2 \rangle\), with the
hyperbolic-plane summand
\(\HypPlan = \begin{pmatrix} 0 & -1 \\ -1 & 0 \end{pmatrix}\) carrying
signature \((1, 1)\) and the rank-1 summand carrying form value \(+2\)
on a generator. The two writings are the same lattice expressed in
opposite sign conventions on the bilinear form: the convention
\((p, q) = (3, 2)\) is signature ``three positive, two negative''
for the fixed quadratic form, while the
Dolgachev convention writes the same signature ``two positive, three
negative'' on the form with reversed sign. The
\((+2)\)-vectors in the fixed convention are the \((-2)\)-vectors of the
Dolgachev convention; the Heegner divisors and the period domain
are independent of the sign choice.

In the Dolgachev sign on \(\widetilde M\), the period domain attached
to \(\widetilde M\) is the type-IV symmetric domain
\[
\Omega_{\widetilde M}
\;=\;
\bigl\{
[Z] \in \P\!\bigl(\widetilde M \otimes \C\bigr)
\;:\;
(Z, Z) = 0,\;
(Z, \bar Z) > 0
\bigr\}^+,
\qquad
\dim_{\C} \Omega_{\widetilde M} = 3,
\]
whereas the exterior-square sign of \(\Lambda^{3,2}\) writes the same
domain with \((Z,\bar Z)<0\).  The two inequalities define the same
oriented two-plane after multiplying the bilinear form by \(-1\).
The arithmetic group is
\[
\Gamma_{\widetilde M}
:=\wedge^2\Sp_4(\Z)/\{\pm\Id_4\}
\subset \mathrm P\mathrm O^+(\widetilde M)
\]
acting properly discontinuously, where the component is fixed by the
positive cone \(\mathcal C(\Lambda^{2,1})_+\)
(\S\,\ref{ssec:exceptional-iso}). The exceptional isogeny
\[
\PSp_4(\R)
\;\simeq\;
\SO(3, 2)_+,
\qquad
\Sp_4(\Z) / \{\pm 1\}
\;\simeq\;
\Gamma_{\widetilde M},
\]
imported as View~\textcircled{5} (the \(\PSp_4\simeq\SO(3,2)_+\)
bridge), gives the identification
\begin{equation}
\label{eq:view4-period-bridge}
\Omega_{\widetilde M} \big/ \Gamma_{\widetilde M}
\;\;\simeq\;\;
\mathcal H_2 \big/ \Sp_4(\Z).
\end{equation}
The form
\(\Delta_5 \in M_5(\Sp_4(\Z), \nu_{\Delta_5})\) of
View~\textcircled{1}, viewed as a section on the right-hand side of
\eqref{eq:view4-period-bridge}, is the Borcherds lift of the weak Jacobi
form \(\phi_{0,1}\) under the data
\((\widetilde M, \phi_{0,1})\) on the left-hand side
(Borcherds~\cite[Thm.~13.3]{B} applied to signature
\((2, 3)\); Gritsenko--Nikulin~\cite[Thm.~1.1]{GNII} for the
identification with the additive Gritsenko lift; the rank-3 row of
the Vol~III lattice-polarised \(\mathfrak g_L\) catalogue at
\(L = U \oplus \langle -2 \rangle\)).

The K3-side period domain on the rank-3 polarisation is a different,
larger object:
\[
\Omega_M
\;=\;
\bigl\{
[\omega] \in \P\!\bigl(M^\perp \otimes \C\bigr)
\;:\;
\langle \omega, \omega \rangle = 0,\;
\langle \omega, \bar\omega \rangle > 0
\bigr\}^+,
\qquad
\dim_{\C} \Omega_M = 17,
\]
where \(M^\perp \subset \Lambda_{K3}\) has signature \((2, 17)\) and
parametrises the \emph{transcendental} Hodge filtration of a
lattice-polarised K3 surface. The Borcherds lift used in
View~\textcircled{4} is on \(\Omega_{\widetilde M}\), not on
\(\Omega_M\): the difference is the adjoining of one Mukai hyperbolic
plane to the rank-\(3\) K3 lattice, not an identification of the
even elliptic cohomology with the elliptic period.  The two period spaces are not
interchangeable: \(\Omega_{\widetilde M}\) of dimension \(3\) is the
genus-two Siegel space and supports \(\Delta_5\) directly via the
exceptional isogeny, while \(\Omega_M\) of dimension \(17\) is the
K3-only period space, on which the rank-\(3\) lattice-polarised lift
of \(\phi_{0,1}\) lands only after the additional Mukai-\(U\)
augmentation.

\subsection{The K3 \texorpdfstring{$\times$}{x} E Mukai--Hodge quotient}
\label{subsec:k3xe-period-family}

The full Calabi--Yau-threefold variation of Hodge structure on
\(X=S\times E_\tau\) is
\[
H^3(X,\Z)
\;=\;
H^2(S,\Z)\otimes H^1(E_\tau,\Z),
\]
of rank \(44\).  The elliptic modulus \(\tau\) belongs to
\(H^1(E_\tau)\).  It does not belong to the even elliptic lattice
\(H^0(E,\Z)\oplus H^2(E,\Z)\), whose Hodge type is constant.

The Mukai lattice
\(\widetilde H(S,\Z)=H^*(S, \Z)\) of signature \((4, 20)\) carries the Mukai
pairing
\(\langle (r, c, s), (r', c', s') \rangle_{\mathrm{Muk}} = c \cdot c' - r s' - r' s\)
(Definition~\ref{def:mukai-lattice}).  The even cohomology of the elliptic curve
is the hyperbolic plane
\(U_E = H^{\mathrm{even}}(E, \Z) = H^0(E,\Z)\oplus H^2(E,\Z)\)
with the standard hyperbolic form. Since \(H^{\mathrm{odd}}(S)=0\), the even
integral cohomology is, as an integral group with its Mukai tensor pairing,
\[
H^{\mathrm{even}}(K3 \times E, \Z)
\;=\;
\widetilde H(S,\Z) \otimes U_E
\;=\;
\widetilde H(S,\Z) \otimes U,
\]
a lattice of rank \(48\).  This even Mukai lattice is not the
Calabi--Yau-threefold period lattice; it is the lattice from which the
rank-\(3\) Mukai-invariant quotient is cut.

For \((K3 \times E)\)-pairs \((S, E_\tau)\) with \(S\)
lattice-polarised by \(M = U \oplus \langle -2 \rangle\), the retained
rank-\(3\) Mukai--Hodge quotient is
\begin{equation}
\label{eq:view4-k3e-moduli}
\mathcal P_{K3 \times E}^{(3)}
\;\;:=\;\;
\Omega_{\widetilde M} \big/ \Gamma_{\widetilde M}
\;\;\overset{\eqref{eq:view4-period-bridge}}{\simeq}\;\;
\mathcal H_2 \big/ \Sp_4(\Z),
\end{equation}
where the augmented lattice \(\widetilde M=U_{\mathrm{mir}}\oplus M\)
is a rank-\(5\) type-IV lattice.  The summand \(U_{\mathrm{mir}}\) is
an abstract Mukai hyperbolic plane in the period quotient.  It is not
the constant even elliptic lattice \(U_E\).  The elliptic period enters
through \(H^1(E_\tau)\) and then through the exterior-square
\(\PSp_4\simeq\SO(3,2)_+\) coordinate system on
\(\Omega_{\widetilde M}\).
Equation~\eqref{eq:view4-k3e-moduli} is the period quotient on which
View~\textcircled{4} is stated.  The quotient is the
\(\PSp_4\simeq\SO(3,2)_+\) model; the assertion that it is the base of
the compact mirror discriminant is Conjecture~\ref{conj:mirror-discriminant}.

The K3-only period domain \(\Omega_M\) of dimension \(17\) does not
enter \eqref{eq:view4-k3e-moduli}: its role is to parametrise the
transcendental Hodge filtration of \(S\) alone. The rank-\(3\)
\((K3 \times E)\) Mukai--Hodge quotient is on
\(\Omega_{\widetilde M}\), of dimension \(3\). The relation
\(\widetilde M = U_{\mathrm{mir}} \oplus M\) is the rank-five
type-IV bridge, not an equality with the full \(H^3(K3\times E)\)
period domain.

\subsection{The discriminant locus}
\label{subsec:view4-discriminant}

The candidate period quotient carries a natural Heegner divisor.  It
becomes the support of the discriminant divisor of a mirror period
family only after the comparison asserted in
Conjecture~\ref{conj:mirror-discriminant}.
On the type-IV symmetric domain \(\Omega_{\widetilde M}\), the Heegner
divisors are
\[
\mathfrak H_{(\delta)}
\;=\;
\bigl\{
[Z] \in \Omega_{\widetilde M}
\;:\;
(Z, \delta) = 0
\bigr\},
\qquad
\delta \in \widetilde M,
\;\;
(\delta, \delta) = 2,
\]
indexed by primitive square-\(2\) vectors of the fixed form on
\(\widetilde M = \Lambda^{3,2}\)
(equivalently, square-\((-2)\) vectors of the Dolgachev form after the
sign flip recorded above).  Such a hyperplane records an extra
algebraic Mukai vector; on a K3 slice its \(H^2\)-projection is the
ordinary nodal \((-2)\)-class when the projection is nonzero.  The
type-II Heegner divisor is
\begin{equation}
\label{eq:view4-discriminant-divisor}
\mathfrak H_{\mathrm{II}}
\;\;=\;\;
\bigcup_{\substack{\delta \in \widetilde M \\ (\delta, \delta) = 2\\
(\delta,\widetilde M)=2\Z}}
\mathfrak H_{(\delta)} \big/ \Gamma_{\widetilde M}
\;\;\subset\;\;
\Omega_{\widetilde M} / \Gamma_{\widetilde M}.
\end{equation}
The \(\Gamma_{\widetilde M}\)-orbit structure on primitive square-\(2\)
vectors of \(\widetilde M\) breaks into type-I and type-II classes
(\(\delta \cdot \widetilde M = \Z\) versus \(\delta \cdot \widetilde M = 2\Z\),
the dichotomy already imported as the type-II versus type-I
distinction at the \(W^{(2)}(\Lambda_{II}^{2,1})\)-orbit level); the union in
\eqref{eq:view4-discriminant-divisor} runs over the type-II orbit,
which is the orbit relevant to the \(\Sp_4(\Z)\)-image of the
exceptional isogeny. The image in the quotient is one irreducible
Heegner-type divisor
(Bruinier--Funke~\cite[\S\,2]{BruinierFunkeHeegner} for the general
lattice-Heegner formulation in signature \((2, 3)\); the index-1 case
is implicit in Borcherds~\cite[Thm.~10.1]{B}). The same divisor,
transported across the bridge~\eqref{eq:view4-period-bridge}, is the
\emph{Humbert surface}
\[
\mathcal H_2^{(\mathrm{Hum})}
\;\;=\;\;
\bigl\{
[Z] \in \mathcal H_2 \big/ \Sp_4(\Z)
\;:\;
Z = \begin{pmatrix} \tau_1 & 0 \\ 0 & \tau_2 \end{pmatrix} \;\text{up to } \Sp_4(\Z)
\bigr\},
\]
the diagonal locus on which the principally polarised abelian surface
\(\C^2 / (\Z^2 + Z \Z^2)\) decomposes as a product of two elliptic
curves.  Under the lattice identification
\eqref{eq:view4-period-bridge}, this Humbert surface is the Heegner
support of a norm-\((-2)\) class in the rank-\(3\) Mukai-invariant
period domain.  The assertion that it is the full discriminant support
of a compact mirror family is part of Conjecture~\ref{conj:mirror-discriminant}.

The form \(\Delta_5\) of View~\textcircled{1} vanishes precisely along
this divisor: \(\Delta_5 \in M_5(\Sp_4(\Z), \nu_{\Delta_5})\) has a
simple zero on \(\mathcal H_2^{(\mathrm{Hum})}\)
\cite[Thm.~3.1]{GNPartI}\cite[\S\,3]{GNII}, with multiplier
\(\nu_{\Delta_5}\) accounting for the half-integral shift. Concretely,
the Borcherds product expansion of View~\textcircled{1} gives
\[
\Delta_5(Z)
\;=\;
64\, q^{1/2} r^{1/2} s^{1/2}
\prod_{(n, l, m) \in \Gamma_{\mathrm{eff}}}
(1 - q^n r^l s^m)^{f(nm, l)},
\]
and the automorphic zero divisor is exactly
\(\mathcal H_2^{(\mathrm{Hum})}\), counted with multiplicity \(1\).
Thus the theorem-level divisor identity is
\begin{equation}
\label{eq:view4-divisor-of-Delta5}
[\,\operatorname{div}(\Delta_5)\,]
\;\;=\;\;
[\,\mathcal H_2^{(\mathrm{Hum})}\,]
\quad \text{(divisor classes on } \mathcal H_2 / \Sp_4(\Z)).
\end{equation}
The additional equality with the mirror discriminant
\(\mathfrak D_{\mathrm{mir}}\) is not contained in the automorphic
theorem; it is the support part of the mirror-discriminant conjecture.

\subsection{The conjecture}
\label{subsec:view4-conjecture}

The discriminant of a period quotient carries multiplicities, not just
a support.  Let
\(\mathfrak D_{\mathrm{mir}}\) denote the conjectural reduced
discriminant divisor of that quotient, and let
\(\lambda_{\mathrm{per}}=\mathcal F^2\) denote the descended K3
period Hodge line on the same quotient.  In the rank-\(3\)
Mukai-invariant component the expected section of
\(\lambda_{\mathrm{per}}^{-10}\) has polar order \(2\) along the type-II Heegner divisor:
locally this is the Picard--Lefschetz contribution of a nodal K3
vanishing cycle to the Hodge filtration, and automorphically it is the
square of the simple zero of \(\Delta_5\).  This multiplicity comparison
is part of the conjecture until the rank-\(3\) quotient and the
comparison line have been constructed.

\begin{conjecture}[The mirror discriminant identity]
\label{conj:mirror-discriminant}
\ClaimStatusConjectured
On the rank-\(3\) \((K3 \times E)\) Mukai--Hodge quotient
parametrised by \eqref{eq:view4-k3e-moduli}, the meromorphic section
\[
\Delta_5^{\,-2}
\;\in\;
\Gamma\!\left(
\mathcal H_2 / \Sp_4(\Z),
\;\;
\mathcal L^{-10}_{\,\mathrm{aut}}
\right)
\]
of the inverse \(10\)th tensor power of the automorphic line bundle
\(\mathcal L_{\mathrm{aut}}\) coincides, as a meromorphic section of
\(\lambda_{\mathrm{per}}^{-10}\) on
\(\mathcal P_{K3 \times E}^{(3)}\), with the
discriminant section of the rank-\(3\) quotient. Equivalently, the polar
divisor of \(\Delta_5^{-2}\) is twice the quotient discriminant
\(\mathfrak D_{\mathrm{mir}}\), the support of
\(\mathfrak D_{\mathrm{mir}}\) is \(\mathfrak H_{\mathrm{II}}\), and the comparison map
\[
\mathcal L^{-10}_{\,\mathrm{aut}}
\;\;\xrightarrow{\;\;\sim\;\;}\;\;
\lambda_{\mathrm{per}}^{-10}
\]
on the open complement \(\mathcal P_{K3 \times E}^{\circ}
= \mathcal P_{K3 \times E}^{(3)} \setminus \mathfrak D_{\mathrm{mir}}\)
extends to a global identification with \(\Delta_5^{-2}\) as a
section of \(\lambda_{\mathrm{per}}^{-10}\).
\end{conjecture}

The conjecture is, at the level of divisor classes,
\begin{equation}
\label{eq:view4-divisor-equation}
[\,\Delta_5^{\,-2}\,]
\;\;=\;\;
2\,[\,\mathfrak D_{\mathrm{mir}}\,]
\;\;=\;\;
2\,[\,\mathfrak H_{\mathrm{II}}\,]
\;\;=\;\;
2\,[\,\mathcal H_2^{(\mathrm{Hum})}\,];
\end{equation}
the identity~\eqref{eq:view4-divisor-of-Delta5} fixes the right-hand
side, and the conjectural content is the matching of multiplicities
and the trivialisation of the comparison line bundle.

\paragraph{Evidence.}
The genus-two Bryan--Pandharipande generating function for
disconnected genus-two BPS states on \(K3 \times E\) predicts the cusp
branch governed by the reciprocal Igusa form~\cite[Conj.~1]{BRYAN};
Oberdieck--Pixton prove the normalized OP/DT/PT scalar representative
\[
Z^X_{\mathrm{OP}}=-4096\,\Delta_5^{-2}
\]
\cite[Thm.~1]{OB}.  In the trace convention of this manuscript
\[
Z_{\mathrm{BPS}}^{K3\times E}:=-4096^{-1}Z^X_{\mathrm{OP}}
=\Delta_5^{-2}.
\]
This is the input to View~\textcircled{1} on the curve-counting side.  The local
Picard--Lefschetz model and Saito's Hodge-filtration formalism predict
the same order-two degeneration as the automorphic divisor
\([\Delta_5^{-2}]\).  The unproved point is the global comparison: the
line
\(\mathcal L^{-10}_{\,\mathrm{aut}} \otimes \lambda_{\mathrm{per}}^{10}\)
must be trivialised and the identification of the two sides must extend
across the discriminant divisor.

A separate piece of evidence is the rank-\(4\) precedent set by
Gritsenko--Nikulin~\cite[Thm.~4.1]{GNII} in the lattice-polarised
setting at the polarisation \(L\) producing a Borcherds product
identified with the quasi-pull-back of the discriminant section of
the universal lattice-polarised Hodge bundle on
\(\Omega_L / \Gamma_L\). Conjecture~\ref{conj:mirror-discriminant} is
the rank-\(3\) analogue of the same identification, evaluated at the
type-II rank-\(3\) Mukai sublattice
\(M = U \oplus \langle -2 \rangle\) of \(\Lambda^{2,1}\)
(see Appendix~\ref{appendix:mukai-gram-dictionary}); the Fake-Monster Borcherds
product \(\Phi_{12}\) on the Leech-extension lattice
\(\mathrm{II}_{25,1}\) of \cite{BorcherdsGKM88, Borcherds95} is a
distinct, signature-\((26,2)\) object and is not invoked here.

\paragraph{Epistemic strength.}
\ClaimStatusConjectured\ The identification claim of
Conjecture~\ref{conj:mirror-discriminant} is conjectured. The
comparison line bundle is a standard Hodge-theoretic object; the
automorphic Heegner identity~\eqref{eq:view4-divisor-of-Delta5} is a
theorem; the equality
\(\operatorname{Supp}(\mathfrak D_{\mathrm{mir}})=\mathfrak H_{\mathrm{II}}\),
the multiplicity \(2\), and the global extension across the
discriminant divisor are conjectural.

\paragraph{Consequences.}
If proved, View~\textcircled{4} embeds the \(\Delta_5\) identity inside
mirror symmetry on \(K3 \times E\) at the level of period-side
discriminants.  The identification would give a period-side derivation
of \(\Delta_5\) independent of the Borcherds product
(View~\textcircled{1}), the BKM denominator (View~\textcircled{2}),
the Pfaffian of the Dirac--Igusa pro-object (View~\textcircled{3}), and
the singular theta lift (View~\textcircled{5}).  It is a conjectural
fifth chain.

A proof of Conjecture~\ref{conj:mirror-discriminant} would identify
\(\Delta_5^{-2}\) with the section of \(\lambda_{\mathrm{per}}^{-10}\)
on the rank-\(3\)
\((K3 \times E)\) Mukai--Hodge quotient; together with the Bryan--Oberdieck--Pixton
theorem~\cite[Thm.~1]{OB} that
\(Z^X_{\mathrm{OP}}=-4096\,\Delta_5^{-2}\), this would give a
period-side interpretation of the unscaled trace
\(Z_{\mathrm{BPS}}^{K3\times E}:=-4096^{-1}Z^X_{\mathrm{OP}}\).

\subsection{Range of the period conjecture}
\label{subsec:view4-not-asserted}

The Beilinson cut is in force: a statement is not primitive if it
holds only after passing to a trace, choosing a chamber, completing a
category, or installing descent data.
The conjectural identification is on the period side: it concerns
the section \(\Delta_5^{-2}\) of an automorphic line bundle and the
discriminant divisor of the rank-\(3\) \(K3\times E\) mirror period
component. It is not a
claim about the protected operator-level datum
\(\mathfrak D_X^{\mathrm{DI}}\) of View~\textcircled{3}; that datum
is established at the Pfaffian and orientation level only relative to
the retained data in the
View~\textcircled{2}\(\leftrightarrow\)\textcircled{3} chapters and
Appendix~\ref{app:obstruction-discharges}. View~\textcircled{4} does
not enter the operator chain.

In the trace convention of Chapter~\ref{ch:zbps-realization}, the
level-\(n\) scalar partition function
\(Z^{K3 \times E}_{\mathrm{BPS}}:=\Delta_5^{-2}\) is the
orientation-forgetful trace of the Pfaffian datum;
its three-dimensional gravity-line realization requires the
Hall--Borcherds residual (Vol~II, \cite{ChiralBarCobarVolII}). The
mirror conjecture above says nothing about the gravitational lift; it
addresses only the geometric origin of the section
\(\Delta_5^{-2}\) on the rank-\(3\) mirror period component.

The Borcherds--Kac--Moody denominator
\(\operatorname{den}(\mathfrak g_{\Delta_5}) = 64^{-1} \Delta_5(2 Z)\)
is on the algebraic side; it is a theorem of
Gritsenko--Nikulin~\cite[\S\,3, Thm.~3.1]{GNII}. The mirror
conjecture is on the period side and does not enter the algebraic
chain.

\subsection{Lattice-convention check}
\label{subsec:view4-lattice-check}

The lattice convention adopted above is the Dolgachev convention
\(\Lambda_{K3} = 3 U \oplus 2 E_8(-1)\), \(M = U \oplus \langle -2 \rangle\),
mirror lattice
\(\check M \simeq U \oplus E_7(-1) \oplus E_8(-1)\) (rank \(17\),
signature \((1,16)\)). Aspinwall--Morrison's original mirror-pair
construction~\cite{AspinwallMorrisonStringy} adopts the same lattice up
to the choice of which hyperbolic plane in \(\Lambda_{K3}\) is split
off; the two conventions agree on the K3 mirror lattice. The
Mukai-lattice form
\(\widetilde M = U \oplus M = 2U \oplus \langle -2\rangle\)
of rank 5 and signature \((2, 3)\) is the convention adopted in
Vol~III~\cite[Prop.~6.4]{CYQGVolIII}; it is the convention
relevant to the period-domain bridge~\eqref{eq:view4-period-bridge}.
The lattice \(\widetilde M\) is identified with
\(\Lambda^{3,2}\) of \S\,\ref{ssec:exceptional-iso} by the sign-flip
on the rank-\(1\) summand recorded in \S\,\ref{subsec:k3-mirror-period}.
Disagreement on the lattice convention appears
only as a sign convention on the Mukai pairing; it does not affect the
discriminant divisor.

The K3 vs.\ \((K3 \times E)\) period-space matching
\eqref{eq:view4-k3e-moduli} adjoins the distinguished hyperbolic
summand \(U_{\mathrm{mir}}\) to the rank-\(3\) K3 lattice and lets the
elliptic variation enter through \(H^1(E_\tau)\).  It does not
identify \(U_{\mathrm{mir}}\) with the constant even elliptic lattice
\(U_E\), and it is not an abstract equality between the
K3-mirror period family and the \((K3 \times E)\)-mirror period family.
The possible discrepancy is exactly the comparison between the
\(\Sp_4\)-side and the \(\mathrm O^+(2,3)\)-side; the
isomorphism~\eqref{eq:view4-period-bridge} is
exactly the \(\PSp_4\simeq\SO(3,2)_+\) bridge of
View~\textcircled{5}.

The word discriminant names a section, not only a support, once
multiplicities are present. There are two natural divisors on the
period side: the Heegner support \(\mathfrak H_{\mathrm{II}}\), and
the polar divisor of the section of \(\lambda_{\mathrm{per}}^{-10}\).
The conjecture concerns the second:
\(\Delta_5^{-2}\) is conjectured to be that section, with polar
divisor twice the support divisor \(\mathfrak H_{\mathrm{II}}\).  The
factor of \(2\) is predicted by the
local Picard--Lefschetz model and is matched, on the automorphic side,
by the square of \(\Delta_5\); a mismatch at the level of multiplicity
would falsify the conjecture by a factor that the Borcherds product
does not allow.

The relation to the Borcherds--Gritsenko--Nikulin quasi-pull-back
theorems is the cleanest comparison.  There the Borcherds form
\(\Phi_{12}\) on the even unimodular lattice of signature \((2,26)\)
is pulled back to lattice-polarised type-IV domains, and its divisor is
read as a discriminant section of the corresponding Hodge bundle
\cite[Thm.~4.1]{GNII}.  The rank-\(3\) lattice
\(L = U \oplus \langle -2 \rangle\) requires a separate comparison line
and a separate Heegner-divisor analysis; the \(\Phi_{12}\) theorem is
evidence, not a proof.

\subsection{Components of the period conjecture}
\label{subsec:view4-period-components}

\begin{enumerate}
\item[\textnormal{(M1)}]
\ClaimStatusOpen\ \emph{Trivialisation of the comparison line.} The
comparison line bundle
\(\mathcal L^{-10}_{\,\mathrm{aut}} \otimes \lambda_{\mathrm{per}}^{10}\)
on \(\mathcal P_{K3 \times E}^{\circ}\) must be trivialised.  Since
\(\nu_{\Delta_5}^2=1\), the inverse-square section has no residual
character; the obstruction is an ordinary Hodge line class.  A
trivialisation of this line, together with the divisor and boundary
conditions below, strengthens the divisor comparison to an equality of
meromorphic sections.

\item[\textnormal{(M2)}]
\ClaimStatusOpen\ \emph{Extension and multiplicity.}
The section comparison is first defined on
\(\mathcal P_{K3 \times E}^{\circ}\), where the Heegner divisor has
been removed.  The period condition is the extension across
\(\mathfrak H_{\mathrm{II}}\) with polar order \(2\), together with
the extension to the rank-2 cusp with boundary exponent governed by
the Borcherds singular weight \(c_1(0)/2=5\)
(\cite[\S\,15]{Borcherds95}).  Thus the divisor of the extended
section, not the open part itself, is compared with
\([\Delta_5^{-2}]\).

\item[\textnormal{(M3)}]
\ClaimStatusOpen\ \emph{The mirror correspondence on the section.}
Dolgachev mirror symmetry gives a correspondence between the
\(M\)-polarised and \(\check M\)-polarised period quotients, not a
canonical involution of \(\Omega_{\widetilde M}/\Gamma_{\widetilde M}\).
The period condition is that a chosen correspondence identifies the
period Hodge line with the automorphic line and carries the
type-II Heegner divisor to the type-II Heegner divisor.  After this
comparison, \(\Delta_5^{-2}\) is required to be the same meromorphic
section of the compared inverse Hodge line.

\item[\textnormal{(M4)}]
\ClaimStatusOpen\ \emph{Rank \(\geq 5\) lattice compatibility.}
The Vol~III rank-\(\geq 5\) lattice-polarised \(\mathfrak g_L\)
family~\cite{CYQGVolIIILatticePolarized} is the conjectural extension of
Conjecture~\ref{conj:mirror-discriminant} to higher Picard rank.
At rank \(\geq 5\) the discriminant identity is conditional on the
Gritsenko--Cl\'ery 2018 admissibility conjecture
(Conj.~5.1 of \cite{GC}); the rank-\(3\) statement is the corresponding
conjectural statement for the K3-Mukai-invariant lattice and is the
rank relevant to the \(\Delta_5\) row.
\end{enumerate}

The four period conditions \((\mathrm M 1)\)--\((\mathrm M 4)\) are
distinct from the four obstruction classes \((\mathrm{D0})\),
\((\mathrm{O1})\), \((\mathrm{O1}^+)\), \((\mathrm{O2})\) of
View~\textcircled{3}; the latter govern the operator-level chain
and are the conditioning for \textcircled{2}\(\leftrightarrow\)\textcircled{3},
while the former govern the period-side chain and are the
conditioning for \textcircled{4}\(\leftrightarrow\)\textcircled{1}.
The two systems lie on different chains: \((\mathrm M 1)\)--\((\mathrm M 4)\)
are period-line and Heegner-divisor conditions, whereas
\((\mathrm{D0})\), \((\mathrm{O1})\), \((\mathrm{O1}^+)\), and
\((\mathrm{O2})\) are Pfaffian-line and orientation conditions.

\subsection{Compatibility with the Vol III lattice convention}
\label{subsec:view4-volIII-lattice-convention}

The lattice convention
\(M = U \oplus \langle -2 \rangle\),
\(\check M \simeq U \oplus E_7(-1) \oplus E_8(-1)\),
\(\widetilde M = U \oplus M\) of
\S\,\ref{subsec:k3-mirror-period} is the convention of
Vol~III~\cite[\S\,6.4 + Prop.~6.4 + Thm.~6.7]{CYQGVolIII}; on the
\((K3 \times E)\)-component the Vol~III \(\kappa\)-tuple
\(\bigl(\kappa_{\mathrm{cat}}, \kappa_{\mathrm{ch}}^{\mathrm{Hodge}},
\kappa_{\mathrm{ch}}^{\mathrm{Heis}}, \kappa_{\mathrm{BKM}}, \kappa_{\mathrm{fiber}}\bigr)
= (0, 0, 3, 5, 24)\) is recorded with each entry sourced from a
distinct construction. The mirror discriminant of
Conjecture~\ref{conj:mirror-discriminant} is on the
\(\kappa_{\mathrm{BKM}} = 5\) row of the \(\kappa\)-tuple; the
period-side datum is the discriminant divisor on the
lattice-polarised period stratum of the Vol~III chart
classification (Vol~III~\cite[\S\,V]{CYQGVolIII}: the
``lattice-polarised period domain'' equivariance stratum). The
\((K3 \times E)\)-fiber rank \(\kappa_{\mathrm{fiber}} = 24\) is the
Mukai-lattice rank of \(\widetilde \Lambda_{K3}\); it is the rank of
the Heisenberg algebra of the Mukai lattice, not a discriminant
invariant.

The mixed-holomorphic-topological-strings companion volume
\cite{MixedHTSCompanion} records that the modular line bundle on
\(\overline{\mathcal A}_g\) at level \(g = 2\) is
\(\Omega_{\mathrm{central}}\), and that
\(\Phi_{10}^{\mathrm{un}} = \chi_{10} = \Delta_5^{\,2}\) is the genus-\(2\)
section.  The mirror-discriminant conjecture is consistent with that statement:
the polar divisor of \(\Delta_5^{-2} = (\Phi_{10}^{\mathrm{un}})^{-1}\) on
\(\overline{\mathcal A}_2\) is twice the divisor of \(\Delta_5\),
and the mirror discriminant identification is precisely the
geometric origin of that section.

The Vol~II Hall--Borcherds residual is silent on
View~\textcircled{4}: the period-side discriminant identification
is independent of the gravitational-lift bridge, and the
conjecture above does not require the residual.
Solving the Hall--Borcherds residual on the
\((K3 \times E)\)-component would not by itself prove
Conjecture~\ref{conj:mirror-discriminant}; conversely, a proof
of Conjecture~\ref{conj:mirror-discriminant} would not
automatically construct the Hall--Borcherds residual.

\medskip

View~\textcircled{4} remains conjectural.  The automorphic, BKM,
Pfaffian, and theta-lift views form the theorem-level quadrilateral at
the stated categorical dimensions.  The mirror-discriminant identity is
the remaining period-side equality.
